\documentclass[border=8pt]{standalone}
\usepackage{lmodern}
\usepackage[T1]{fontenc}
\usepackage{amsmath,amssymb}
\usepackage{tikz}
\begin{document}
% A paragraph set inside a circle. \parshape takes one (indent, width) pair
% per line; pgfmath computes each width as a chord of the circle.
\def\R{130.9}               % radius in pt (4.6cm)
\def\B{12}                  % \baselineskip in pt
\newcount\n \n=20           % lines: the text block is n\B = 240pt tall
\newcount\i
\def\shape{}
\i=1
\loop
  \pgfmathsetmacro\yy{(\n*\B/2-(\i-0.5)*\B)/10}    % line centre, relative to circle centre, in 10pt units
  \pgfmathsetmacro\hw{10*sqrt(\R/10*\R/10-\yy*\yy)-6} % half chord, 6pt inside the rim (pgfmath overflows above 16383, hence the /10)
  \pgfmathsetmacro\ind{\R-\hw}
  \pgfmathsetmacro\lw{2*\hw}
  \edef\shape{\shape\ind pt \lw pt }
  \advance\i 1
\ifnum\i<\numexpr\n+1\relax\repeat
\hbox{%
  \rlap{\tikz[baseline=0pt]{\fill[blue!6] (\R pt,-111.6pt) circle (\R pt); \draw[blue!50!black,line width=.5pt] (\R pt,-111.6pt) circle (\R pt);}}%
  \vtop{\hsize=2\dimexpr\R pt\relax \baselineskip=\B pt \parindent=0pt \tolerance=3000 \emergencystretch=1.5em \hyphenpenalty=10
    \parshape \n \shape
    The ratio of a circle's circumference to its diameter,
    $\pi\approx3.14159\ldots$, is the same for every circle
    \tikz[baseline=-0.6ex]{\draw (0,0) circle (0.8ex); \draw (-0.8ex,0)--(0.8ex,0);}.
    Archimedes squeezed it between inscribed and circumscribed polygons
    \tikz[baseline=-0.6ex]{\draw (0,0) circle (0.8ex); \draw (0:0.8ex) \foreach \a in {60,120,...,300} {-- (\a:0.8ex)} -- cycle;},
    proving $3\tfrac{10}{71}<\pi<3\tfrac17$; two thousand years later Leibniz found
    \[ \frac{\pi}{4}=\sum_{k=0}^{\infty}\frac{(-1)^k}{2k+1}=1-\frac13+\frac15-\frac17+\cdots, \]
    and Euler tied $\pi$ to the primes through $\sum_{n\ge1}n^{-2}=\pi^2/6$
    and to $e$ through $e^{i\pi}+1=0$; Wallis wrote $\frac{\pi}{2}=\prod_{n\ge1}\frac{4n^2}{4n^2-1}$.
    The Gaussian integral
    $\int_{-\infty}^{\infty}e^{-x^2}\,dx=\sqrt{\pi}$
    \tikz[baseline=-0.3ex]{\draw[thick] plot[domain=-2.2:2.2,samples=30] ({\x*0.42em},{exp(-\x*\x)*1.4ex});}
    carries it into probability and statistics, and it hides in Stirling's
    formula $n!\sim\sqrt{2\pi n}\,(n/e)^n$. This paragraph is set by
    \TeX's \texttt{\char`\\parshape} primitive: twenty line widths
    computed by pgfmath as chords of the circle, with the
    display and the inline TikZ pictures flowing
    through the same shape as the prose.\par}}
\end{document}
